\documentclass[10pt,a4paper]{article}
\usepackage[utf8]{inputenc}
\usepackage[T1]{fontenc}
\usepackage{graphicx}
\usepackage{booktabs}
\usepackage{float}
\usepackage{listings}
\usepackage[margin=1.7cm]{geometry}
\usepackage{xcolor}
\usepackage{amsmath}
\usepackage{hyperref}

\setlength{\parskip}{0.2ex}
\setlength{\parindent}{0pt}

\lstset{
  basicstyle=\ttfamily\small,
  breaklines=true,
  frame=single,
  backgroundcolor=\color{gray!10},
  numbers=left,
  numberstyle=\tiny\color{gray},
}

\title{COMP90051 Group Project Title}
\author{Student ID -- Name \and Student ID -- Name}
\date{}

\begin{document}
\maketitle

\section{Introduction}

State the research question, dataset(s), public URL(s), and why the question goes
beyond the obvious dataset task.

\section{Literature Review}

Summarise dataset references and at least two related papers. Include the paper
for the recent algorithm selected under the project specification.

\section{Methods}

Describe feature construction, preprocessing, three algorithms, outer
cross-validation, nested cross-validation, and hyperparameter tuning.

\subsection{Feature Construction and Preprocessing}

Describe the non-trivial transformations applied before model fitting. Make clear
which transformations are fit only on training folds to avoid leakage.

\subsection{Algorithms}

Summarise the three algorithms and their model class complexity. For the
2016-or-later research-paper algorithm, cite the qualifying paper and explain the
core idea at a postgraduate-readable level.

\begin{table}[H]
\centering
\small
\begin{tabular}{llll}
\toprule
Complexity & Algorithm & Tuned parameter(s) & Citation \\
\midrule
Simple & TBD & TBD & TBD \\
Medium & TBD & TBD & TBD \\
Complex & TBD & TBD & TBD \\
\bottomrule
\end{tabular}
\caption{Algorithms selected for comparison.}
\end{table}

\subsection{Cross-Validation and Tuning}

Describe the outer cross-validation used for fair comparison and the nested
cross-validation used for hyperparameter selection. State that these were
implemented from scratch, and list the candidate values considered.

\begin{table}[H]
\centering
\small
\begin{tabular}{lll}
\toprule
Algorithm & Hyperparameter & Candidate values \\
\midrule
TBD & TBD & TBD \\
TBD & TBD & TBD \\
TBD & TBD & TBD \\
\bottomrule
\end{tabular}
\caption{Hyperparameter grids used in nested cross-validation.}
\end{table}

\section{Results and Discussion}

Report at least three experimental results with error bars. Discuss why methods
worked or did not work, their advantages and disadvantages, and alternatives
considered.

\begin{table}[H]
\centering
\small
\begin{tabular}{lrrr}
\toprule
Algorithm & Metric 1 & Metric 2 & Metric 3 \\
\midrule
TBD & TBD & TBD & TBD \\
TBD & TBD & TBD & TBD \\
TBD & TBD & TBD & TBD \\
\bottomrule
\end{tabular}
\caption{Experimental results. Report mean values with error bars.}
\end{table}

\section{Conclusion}

Summarise the knowledge gained about the problem and the research question.

\begin{thebibliography}{9}

\bibitem{dataset_placeholder}
Dataset Authors. (2026).
\textit{Dataset reference placeholder}.
\url{https://example.com}.
Replace this entry with the actual dataset citation and add the selected papers.

\end{thebibliography}

\end{document}
